\documentclass[11pt, a4paper]{article}

\usepackage[utf8]{inputenc}
\usepackage[T1]{fontenc}
\usepackage[spanish]{babel}
\usepackage[margin=2cm, headheight=15pt]{geometry} % Se ajustó headheight para fancyhdr
\usepackage{fancyhdr}
\usepackage{amsmath, amssymb}
\usepackage{siunitx}
\usepackage[american, RPvoltages]{circuitikz}

\pagestyle{fancy}
\fancyhf{}
\lhead{\textbf{IMT-121 | Diagnóstico EC1 y Ticket de Salida}}
\rhead{\textbf{Sesión Semana 5}}

\begin{document}

\begin{center}
    {\Large \textbf{Checkpoint Formativo y Ticket de Salida}} \\[0.2cm]
    \textit{Identificación de Modelos Matriciales y Principio de Linealidad}
\end{center}

\textbf{Estudiante:} \rule{8cm}{0.4pt} \hfill \textbf{Fecha:} \rule{3cm}{0.4pt}
\vspace{0.4cm}

\section*{Parte 1: Diagnóstico de Formulación Matricial (EC1)}
Analiza el siguiente circuito en CC y responde seleccionando el modelo matricial $\mathbf{R}\mathbf{I} = \mathbf{V}$ correcto. 

\vspace{0.2cm}
\begin{minipage}{0.45\textwidth}
    \begin{center}
        \begin{circuitikz}[scale=0.8, transform shape]
            \draw (0,0) to[V, v=$V_1$] (0,3)
                  to[R, l=$R_A$, i=$i_1$] (3,3)
                  to[R, l=$R_B$] (3,0) -- (0,0);
            \draw (3,3) to[R, l=$R_C$, i=$i_2$] (6,3)
                  to[V, v_=$V_2$] (6,0) -- (3,0);
            \draw (3,0) node[ground]{};
        \end{circuitikz}
    \end{center}
\end{minipage}
\begin{minipage}{0.5\textwidth}
    \textbf{Selecciona la matriz correcta:}
    
    \vspace{0.2cm}
    $\square$ \textbf{A:} $\begin{bmatrix} (R_A+R_B) & R_B \\ R_B & (R_B+R_C) \end{bmatrix} \begin{bmatrix} i_1 \\ i_2 \end{bmatrix} = \begin{bmatrix} V_1 \\ V_2 \end{bmatrix}$
    
    \vspace{0.2cm}
    $\square$ \textbf{B:} $\begin{bmatrix} (R_A+R_B) & -R_B \\ -R_B & (R_B+R_C) \end{bmatrix} \begin{bmatrix} i_1 \\ i_2 \end{bmatrix} = \begin{bmatrix} V_1 \\ -V_2 \end{bmatrix}$
    
    \vspace{0.2cm}
    $\square$ \textbf{C:} $\begin{bmatrix} R_A & -R_B \\ -R_B & R_C \end{bmatrix} \begin{bmatrix} i_1 \\ i_2 \end{bmatrix} = \begin{bmatrix} V_1 \\ -V_2 \end{bmatrix}$
    
    \vspace{0.2cm}
    $\square$ \textbf{D:} $\begin{bmatrix} (R_A+R_B) & -R_B \\ -R_B & (R_B+R_C) \end{bmatrix} \begin{bmatrix} i_1 \\ i_2 \end{bmatrix} = \begin{bmatrix} -V_1 \\ V_2 \end{bmatrix}$
\end{minipage}

\vspace{0.5cm}
\textbf{Justificación breve:} ¿Por qué el elemento $(2,2)$ de la matriz es $(R_B + R_C)$ y por qué el término independiente de la malla 2 tiene el signo que elegiste? \\
\rule{\textwidth}{0.4pt} \\
\rule{\textwidth}{0.4pt}

\vspace{0.4cm}
\hrule
\vspace{0.4cm}

\section*{Parte 2: Ticket de Salida (Transición a Superposición)}
El siguiente circuito lineal posee dos fuentes independientes ($V_s$ y $I_s$). Observa cómo están conectadas.

\begin{center}
    \begin{circuitikz}[scale=0.9, transform shape]
        \draw (0,0) to[V, v=$V_s$] (0,3)
              to[R, l=$R_1$] (3,3)
              to[R, l=$R_2$, i=$I_x$, color=blue] (3,0) -- (0,0);
        \draw (3,3) to[R, l=$R_3$] (6,3)
              to[I, i_=$I_s$] (6,0) -- (3,0);
        \draw (3,0) node[ground]{};
    \end{circuitikz}
\end{center}

\textbf{1. Variable de Respuesta:} ¿Cuál es la variable eléctrica que estamos evaluando en el componente central? \\
\rule{8cm}{0.4pt}

\vspace{0.3cm}
\textbf{2. Contribuciones:} Por linealidad, la respuesta total es la suma de las contribuciones aisladas. Escribe la ecuación cualitativa que describe esto: \\
$I_{x(\text{total})} =$ \rule{5cm}{0.4pt} $+$ \rule{5cm}{0.4pt}

\vspace{0.3cm}
\textbf{3. Predicción Cualitativa:}
\begin{itemize}
    \item Si ``apagamos'' $I_s$ y solo actúa $V_s$, ¿la corriente aportada fluirá hacia \textbf{abajo} o hacia \textbf{arriba} por $R_2$? \rule{2cm}{0.4pt}
    \item Si ``apagamos'' $V_s$ y solo actúa $I_s$, ¿la corriente aportada fluirá hacia \textbf{abajo} o hacia \textbf{arriba} por $R_2$? \rule{2cm}{0.4pt}
\end{itemize}

\end{document}
